\section{Problem und Ausgangszustand}
\label{chap:problem-initial-state}

Eine frühere Lifecycle-Architektur behandelte ein Full-Stack-Template als Master und
leitete daraus Produktbäume ab. Dieses Modell ist kein aktueller Sollzustand: feste
Wurzeln, Annahmen über Komponenten und Kopplung an eine fremde Git-Historie erschweren
den Austausch eines Toolingordners in bestehenden Projekten.

Der aktuelle Ansatz trennt austauschbaren Toolingbestand von Projektkonfiguration und
State. Ein Austausch von \texttt{tools/} und \texttt{docs/toolingdocs/} wird als
versionierter Migrationsfall behandelt, nicht als gleichversionige Kopie. Abbildung
\ref{fig:project-context} verdeutlicht die abgeleiteten, statt fest codierten Pfade.
\CaseDiagram{project-context-paths}{ProjectContext trennt Projekt-, Tooling-, Dokumentations- und State-Wurzeln.}{fig:project-context}
